\documentclass[../../../main]{subfiles}

\ifSubfilesClassLoaded{%
    \zexternaldocument*{../1.Category-Theory/exported/aux/main}%
    \zexternaldocument*{../1.Additive-and-abelian-categories/exported/aux/main}%
    \zexternaldocument*{../2.Complexes-and-cohomology/exported/aux/main}%
    \zexternaldocument*{../3.Cones-and-homotopy/exported/aux/main}%
    \zexternaldocument*{../4.Projectives-and-injectives/exported/aux/main}%
    \zexternaldocument*{1.Delta-functors}
}{}

\begin{document}
\section{Derived functors}

\begin{definition} \label{def:right-derived-functor}
  Let \(F: \category{A} \to \category{B}\) be a \hyperref[def:exact-complex-functor]{left exact functor} between two \hyperref[def:abelian-category]{abelian categories} \(\category{A}\) and \(\category{B}\).
  If \(\category{A}\) \hyperref[def:enough-injectives-projectives]{has enough injectives} we define the \(i\)-th \textbf{right derived functor} (\(i \geq 0\)) as
  \begin{equation*}
    \forall A \in \category{A} \quad R^{i}F (A) \coloneq \cohomology{i}(F(I^{\bullet}))
  \end{equation*}
  where \(I^{\bullet}\) is an \hyperref[def:injective-projective-resolution]{injective resolution} of \(A\).

  Note that this is well-defined by \autoref{cor:injective-have-same-cohomology}.
\end{definition}

\begin{remark}
  Since
  \begin{center}
    \begin{tikzcd}
      0 \arrow[r] & F(A) \arrow[r] & F(I^{0}) \arrow[r] & F(I^{1})
    \end{tikzcd}
  \end{center}
  is \hyperref[def:short-exact-sequence]{exact} we always have \(R^{0}F(A) \cong F(A)\).
\end{remark}

\begin{lemma}
  Let \(F: \category{A} \to \category{B}\) be a \hyperref[def:exact-complex-functor]{left exact functor} between two \hyperref[def:abelian-category]{abelian categories} \(\category{A}\) and \(\category{B}\).
  For every \(i \in \mathbb{N}\) the \hyperref[def:right-derived-functor]{right derived functor} \(R^{i}F\) is an \hyperref[def:additive-functor]{additive functor} from \(\category{A}\) to \(\category{B}\).
\end{lemma}

\begin{theorem}
  Let \(F: \category{A} \to \category{B}\) be a \hyperref[def:exact-complex-functor]{left exact functor} between two \hyperref[def:abelian-category]{abelian categories} \(\category{A}\) and \(\category{B}\).
  The \hyperref[def:right-derived-functor]{right derived functors} \(R^{\bullet}F\) form a \hyperref[def:delta-functor]{cohomological \(\delta\)-functor}.
\end{theorem}

\begin{theorem}
  Let \(F: \category{A} \to \category{B}\) be a \hyperref[def:exact-complex-functor]{left exact functor} between two \hyperref[def:abelian-category]{abelian categories} \(\category{A}\) and \(\category{B}\).
  Suppose that \(\category{A}\) \hyperref[def:enough-injectives-projectives]{has enough injectives}.
  Then for any object \(A \in \category{A}\) the \hyperref[def:right-derived-functor]{right derived functors} \(R^{\bullet}F\) form an \hyperref[def:universal-delta-functor]{universal \(\delta\)-functor}.
\end{theorem}

\begin{definition} \label{def:acyclic-object}
  Let \(F: \category{A} \to \category{B}\) be a \hyperref[def:exact-complex-functor]{left exact functor} between two \hyperref[def:abelian-category]{abelian categories} \(\category{A}\) and \(\category{B}\).
  An object \(Q \in \category{A}\) is called \(F\)-\textbf{acyclic} if
  \begin{equation*}
    \forall i > 0 \quad R^{i}F(Q) = 0
  \end{equation*}
\end{definition}

\begin{example}
  Every \hyperref[def:injective-projective-object]{injective object} \(I \in \category{A}\) is \hyperref[def:acyclic-object]{\(F\)-acyclic}.
  Indeed, we can compute \(R^{i}F(I)\) using the resolution \(0 \to I \xrightarrow{\id} I \to 0\) which implies that for all \(i > 0 \quad R^{i}F(I) = 0\).
\end{example}


\begin{theorem} \label{thm:dimension-shifting}
  Let \(F: \category{A} \to \category{B}\) be a \hyperref[def:exact-complex-functor]{left exact functor} between two \hyperref[def:abelian-category]{abelian categories} \(\category{A}\) and \(\category{B}\).
  If
  \begin{center}
    \begin{tikzcd}
      0 \arrow[r] & A \arrow[r] & Q^{0} \arrow[r] & Q^{1} \arrow[r] & \cdots \arrow[r] & Q^{m} \arrow[r] & M \arrow[r] & 0
    \end{tikzcd}
  \end{center}
  is \hyperref[def:short-exact-sequence]{exact} with \(Q^{i}\) \hyperref[def:acyclic-object]{\(F\)-acyclic} then for all \(i \geq m + 2 \quad R^{i}F(A) \cong R^{i-m-1}F(A)\) and \(R^{m+1}F(A)\) is the \hyperref[def:categorical-cokernel]{cokernel} of \(F(I^{m}) \to F(M)\).

  In particular if \(A[0] \to Q^{\bullet}\) is an \hyperref[def:acyclic-object]{\(F\)-acyclic resolution} of \(A\) then
  \begin{align*}
    R^{i}F(A) \cong \cohomology{i}(F(Q^{\bullet}))
  \end{align*}
\end{theorem}
\begin{proof}
  We will prove this using induction.

  In the case \(m = 0\) we have the \hyperref[def:short-exact-sequence]{short exact sequence}
  \begin{center}
    \begin{tikzcd}
      0 \arrow[r] & A \arrow[r] & Q^{0} \arrow[r] & M \arrow[r] & 0
    \end{tikzcd}
  \end{center}
  Using the \hyperref[thm:LES-in-cohomology]{long exact sequence in cohomology} we can extend this sequence to
  \begin{center}
    \begin{tikzcd}
      0 \arrow[r] & F(A) \arrow[r] & F(Q^{0}) \arrow[r] \arrow[d, phantom, ""{coordinate, name=X}] & F(M) \arrow[dll, rounded corners, to path={
        -- ([xshift=2ex]\tikztostart.east)
        |- (X) [near end]\tikztonodes
        -| ([xshift=-2ex]\tikztotarget.west)
        -- (\tikztotarget)}, "\delta^{0}" description] \\
      & R^{1}F(A) \arrow[r] & R^{1}F(Q^{0}) \arrow[r] \arrow[d, phantom, ""{coordinate, name=Y}] & R^{1}F(M) \arrow[dll, rounded corners, to path={
        -- ([xshift=2ex]\tikztostart.east)
        |- (Y) [near end]\tikztonodes
        -| ([xshift=-2ex]\tikztotarget.west)
        -- (\tikztotarget)}, "\delta^{1}" description] \\
      & R^{2}F(A) \arrow[r] & R^{2}F(Q^{0}) \arrow[r] & \cdots
    \end{tikzcd}
  \end{center}
  Since \(Q^{0}\) is \(F\)-acyclic we have that \(\forall i > 0 \quad R^{i}F(Q^{0}) = 0\).
  Therefore we have
  \begin{enumerate}
    \item For \(i = 1\) the segment
          \begin{center}
            \begin{tikzcd}
              F(I) \arrow[r] & F(M) \arrow[r, "\delta^{0}"] & R^{1}F(A) \arrow[r] & 0
            \end{tikzcd}
          \end{center}
          By \hyperref[def:short-exact-sequence]{exactness} at \(R^{1}F(A)\) we have
          \begin{equation*}
            R^{1}F(A) \cong \quot{F(M)}{\im F(I)} \cong \coker \left( F(I) \to F(M) \right)
          \end{equation*}
    \item For \(i \geq 2\) the segment
          \begin{center}
            \begin{tikzcd}
              0 \arrow[r] & R^{i-1}F(A) \arrow[r, "\delta^{i-1}"] & R^{i}F(A) \arrow[r] & 0
            \end{tikzcd}
          \end{center}
          By \hyperref[def:short-exact-sequence]{exactness} we have that \(R^{i-1}F(A) \cong R^{i}F(A)\).
  \end{enumerate}
  This finishes the proof of the base case.

  Let's suppose now that we have proven the result for \(m-1\) and check it for \(m\).
  Consider the \hyperref[def:short-exact-sequence]{exact sequence}
  \begin{center}
    \begin{tikzcd}
      0 \arrow[r] & A \arrow[r] & Q^{0} \arrow[r] & Q^{1} \arrow[r] & \cdots \arrow[r] & Q^{m} \arrow[r] & M \arrow[r] & 0
    \end{tikzcd}
  \end{center}
  Denote by \(Z_{0} = A, Z_{k+1} = \coker \left( Z_{k} \to Q^{k} \right)\) and \(Z_{m+1} = M\).
  We can split the original sequence into \(m+1\) \hyperref[def:short-exact-sequence]{short exact sequences}
  \begin{align}
    &0 \to A \to Q^{0} \to Z_{1} \to 0 \label{thm:dimension-shifting-SES-0} \\
    \forall k \in \{1, \ldots, m-1\} \quad & 0 \to Z_{k} \to Q^{k} \to Z_{k+1} \to 0 \label{thm:dimension-shifting-SES-k} \\
    &0 \to Z_{m} \to Q^{m} \to M \to 0 \label{thm:dimension-shifting-SES-m}
  \end{align}
  We now apply the inductive hypothesis to this sequences:
  \begin{enumerate}
    \item Using \eqref{thm:dimension-shifting-SES-0} for all \(i \geq 2\) we get
          \begin{equation*}
            R^{i}F(A) \cong R^{i-1}F(Z_{1})
          \end{equation*}
    \item Using \eqref{thm:dimension-shifting-SES-k} for all \(i \geq m + 2\) we have
          \begin{equation*}
            R^{i}F(A) \cong R^{i-(m+1)}F(Z_{m+1}) \cong R^{i-m-1}F(M)
          \end{equation*}
    \item If \(i = m+1\) we get
          \begin{align*}
            R^{m+1}F(A) \cong R^{1}F(Z_{m}) \cong \coker \left( F(Q^{m}) \to F(M) \right)
          \end{align*}
  \end{enumerate}
  This concludes the proof of the first part of the theorem.

  To see the second part, let
  \begin{center}
    \begin{tikzcd}
      0 \arrow[r] & A \arrow[r] & Q^{0} \arrow[r] & Q^{1} \arrow[r] & \cdots
    \end{tikzcd}
  \end{center}
  be a \hyperref[def:acyclic-object]{\(F\)-acyclic resolution}. We want to see that \(R^{i}F(A) \cong H^{i}(F(Q^{\bullet}))\).

  Fix \(i \geq 1\) and truncate the sequence to form a \hyperref[def:short-exact-sequence]{exact sequence}
  \begin{center}
    \begin{tikzcd}
      0 \arrow[r] & A \arrow[r] & Q^{0} \arrow[r] & \cdots \arrow[r] & Q^{i-1} \arrow[r] & Z_{i} \arrow[r] & 0
    \end{tikzcd}
  \end{center}
  where
  \begin{equation*}
    Z_{i} \coloneq \ker \left( Q^{i} \to Q^{i+1} \right) \cong  \im \left( Q^{i-1} \to Q^{i} \right)
  \end{equation*}

  Since \(F\) is \hyperref[def:exact-complex-functor]{left-exact} applying it to the \hyperref[def:short-exact-sequence]{short exact sequence} \(0 \to Z_{i} \to Q^{i} \to Q^{i+1}\) we get the \hyperref[def:short-exact-sequence]{short exact sequence}
  \begin{center}
    \begin{tikzcd}
      0 \arrow[r] & F(Z_{i}) \arrow[r] & F(Q^{i}) \arrow[r] & F(Q^{i+1})
    \end{tikzcd}
  \end{center}
  Meaning that
  \begin{equation*}
    F(Z_{i}) \cong \ker \left( F(Q^{i}) \to F(Q^{i+1}) \right)
  \end{equation*}
  
  Using the result above (for \(m = i-1\)) we have
  \begin{align*}
    R^{i}F(A) &\cong \coker \left( F(Q^{i-1}) \to F(Z_{i}) \right) \\
              &\cong \quot{ F(Z_{i}) }{\im \left( F(Q^{i-1}) \to F(Z_{i}) \right)} \\
              &\cong \quot{\ker \left( F(Q^{i}) \to F(Q^{i+1}) \right)}{\im \left( F(Q^{i-1}) \to F(Z_{i}) \right)} \\
              &\cong \cohomology{i}(F(Q^{\bullet}))
  \end{align*}
\end{proof}


\end{document}
